% Ad Atticum 13.47 --- headnote
% ad-atticum-13-47, 13 August 45 BC, Tusculum

\textit{Cicero to Atticus, written at the Tusculan villa on 13
August 45 BC --- Perseus dateline \textit{Scr. in Tusculano Id.
Sext. a. 709 (45)}. A single section, vivid and quick. Cicero
addresses Atticus as ``Agamemnon'' --- the great commander
issuing his orders --- and reports, in a quick tricolon, his
obedience: he dropped what he had begun (\textit{instituta
omisi}), threw aside what was in hand (\textit{ea quae in
manibus habebam abieci}), and ``roughed out'' (\textit{edolavi})
what Atticus had told him to write. The \textit{edolavi}
preserves the woodworker's image --- hewn out with the axe, not
yet finished --- and that is precisely the boast: he has done
the rough draft on demand.}

\textit{The rest is housekeeping with one barbed last line. Pollex
will report on the accounts; the freedman is being kept short and
must be looked after this first year, but the leash will be
tightened next. He cannot get down to Puteoli to deal in person
with Cluvius's estate, both because of Torquatus and because
Caesar is now nearby. And then the closing fillip: Dolabella has
written that he is coming the day after the Ides ---
\textit{o magistrum molestum!}, ``what a tiresome schoolmaster!''
Dolabella was taking declamation lessons with Cicero that summer,
and the relentless visits had begun to grate. The whole letter
is one sustained joke about who is master and who is pupil --- of
Agamemnon, of Dolabella, of Cicero himself.}
